% titulo: Preliminares (dedicatoria, siglas, resumen, abstract)
% estado: borrador
% actualizado: 2026-10-19
% La portada se llena en main.tex (\titulo, \autor, ...).

\begin{dedicatoria}
\pendiente{dedicatoria}
\end{dedicatoria}

\begin{agradecimientos}
\pendiente{agradecimientos}
\end{agradecimientos}

\indicegeneral

\begin{siglas}
  \sigla{API}{Application Programming Interface (interfaz de programación de aplicaciones)}
  \sigla{CI}{Continuous Integration (integración continua)}
  \sigla{DLA}{Document Layout Analysis (análisis de la disposición del documento)}
  \sigla{DOI}{Digital Object Identifier}
  \sigla{IEEE}{Institute of Electrical and Electronics Engineers}
  \sigla{IMRyD}{Introducción, Métodos, Resultados y Discusión}
  \sigla{LLM}{Large Language Model (modelo de lenguaje de gran tamaño)}
  \sigla{MVP}{Minimum Viable Product (producto mínimo viable)}
  \sigla{NLP}{Natural Language Processing (procesamiento de lenguaje natural)}
  \sigla{OCR}{Optical Character Recognition (reconocimiento óptico de caracteres)}
  \sigla{OJS}{Open Journal Systems}
  \sigla{PDF}{Portable Document Format}
  \sigla{RE}{Regla editorial (del catálogo de requisitos)}
  \sigla{RICE}{Reach, Impact, Confidence, Effort (método de priorización)}
  \sigla{RNF}{Requisito no funcional}
  \sigla{SDD}{Spec-Driven Development (desarrollo dirigido por especificaciones)}
  \sigla{US}{User Story (historia de usuario)}
  \sigla{YAML}{YAML Ain't Markup Language}
\end{siglas}

\indicetablas

\indicefiguras

\begin{resumen}
\end{resumen}

\begin{abstract}
\pendiente{revisar la traducción}


\textbf{Keywords:} artificial intelligence, Open Journal Systems, editorial review, academic manuscripts, stylistic
suggestions.
\end{abstract}

\begin{palabrasclave}
inteligencia artificial, Open Journal Systems, revisión editorial, manuscritos académicos,
sugerencias estilísticas.
\end{palabrasclave}
